\documentclass[11pt,a4paper]{article}
\usepackage[utf8]{inputenc}
\usepackage[vietnamese]{babel}
\usepackage{amsmath,amssymb}
\usepackage{geometry}
\geometry{a4paper, total={180mm,260mm}, left=15mm, right=15mm, top=18mm, bottom=18mm}
\usepackage{tikz}
\usetikzlibrary{shapes.geometric, arrows.meta, positioning, calc}
\usepackage{tcolorbox}
\tcbuselibrary{skins,breakable}
\usepackage{xcolor}
\usepackage{multicol}
\usepackage{array}
\usepackage{fontawesome}
\usepackage{fancyhdr}

\pagestyle{fancy}
\fancyhf{}
\renewcommand{\headrulewidth}{0.5pt}
\renewcommand{\footrulewidth}{0.5pt}
\fancyhead[L]{\small\bfseries\color{teal}\faBook\ SỔ TAY KIẾN THỨC TOÁN 3}
\fancyhead[R]{\small\bfseries\color{teal}Lớp Toán Thầy Lĩnh -- loptoanthaylinh.vn}
\fancyfoot[C]{\small\color{gray}-- Sổ tay kiến thức: Bài 52. Diện tích hình chữ nhật, diện tích hình vuông --}

\definecolor{primary}{HTML}{005F73}
\definecolor{secondary}{HTML}{0A9396}
\definecolor{accent}{HTML}{EE9B00}
\definecolor{danger}{HTML}{AE2012}
\definecolor{boxbg}{HTML}{F4F6F8}

\begin{document}

% --- BANNER TIÊU ĐỀ ---
\begin{tcolorbox}[
  enhanced,
  colback=white,
  colframe=primary,
  arc=3mm,
  boxrule=1.2pt,
  left=4mm, right=4mm, top=2.5mm, bottom=2.5mm,
  drop shadow=black!10!white
]
  \centering
  {\Large\bfseries\color{primary}\faBookmark\ BÀI 52: DIỆN TÍCH HÌNH CHỮ NHẬT, DIỆN TÍCH HÌNH VUÔNG}\\[1.5mm]
  {\small\itshape\color{black!75}Công thức chuẩn S = a x b và S = a x a; kĩ năng đổi cùng đơn vị và bài toán thực tế trên 1 trang A4}
\end{tcolorbox}

\vspace{-2mm}


\noindent
\begin{minipage}[t]{0.49\textwidth}\vspace{0pt}
\begin{tcolorbox}[
  enhanced,
  colback=blue!4!white,
  colframe=blue!75!black,
  arc=2mm, boxrule=0.8pt,
  left=2mm, right=2mm, top=1.8mm, bottom=1.8mm,
  title={\footnotesize\bfseries\color{white}\faKey\ 1. Diện tích hình chữ nhật}
]
{\footnotesize
\begin{itemize}\setlength{\itemsep}{1.5pt}\setlength{\leftskip}{-2.5mm}
  \item Muốn tính diện tích hình chữ nhật, ta lấy \textbf{chiều dài nhân với chiều rộng} (cùng một đơn vị đo):
  $$\mathbf{S = a \times b}$$
  \item \textbf{Các bài toán ngược:}\\
  Chiều dài $= S : b$ (Diện tích chia cho chiều rộng).\\
  Chiều rộng $= S : a$ (Diện tích chia cho chiều dài).
  \item Đơn vị diện tích tương ứng: $\text{cm} \to \text{cm}^2$; $\text{m} \to \text{m}^2$.
\end{itemize}
}
\end{tcolorbox}
\end{minipage}
\hfill
\begin{minipage}[t]{0.49\textwidth}\vspace{0pt}
\begin{tcolorbox}[
  enhanced,
  colback=teal!4!white,
  colframe=teal!75!black,
  arc=2mm, boxrule=0.8pt,
  left=2mm, right=2mm, top=1.8mm, bottom=1.8mm,
  title={\footnotesize\bfseries\color{white}\faCheckSquare\ 2. Diện tích hình vuông}
]
{\footnotesize
\begin{itemize}\setlength{\itemsep}{1.5pt}\setlength{\leftskip}{-2.5mm}
  \item Muốn tính diện tích hình vuông, ta lấy \textbf{độ dài một cạnh nhân với chính nó}:
  $$\mathbf{S = a \times a}$$
  \item \textbf{Bảng diện tích hình vuông mẫu:}
  \begin{center}
  \renewcommand{\arraystretch}{1.1}
  \begin{tabular}{|c|c|c|c|c|c|}
    \hline
    \textbf{Cạnh ($a$)} & $3\text{ cm}$ & $4\text{ cm}$ & $5\text{ cm}$ & $6\text{ cm}$ & $7\text{ cm}$ \\
    \hline
    \textbf{Diện tích ($S$)} & $9\text{ cm}^2$ & $16\text{ cm}^2$ & $25\text{ cm}^2$ & $36\text{ cm}^2$ & $49\text{ cm}^2$ \\
    \hline
  \end{tabular}
  \end{center}
  \item Muốn tìm cạnh hình vuông khi biết diện tích, ta nhẩm xem số nào nhân với chính nó thì bằng diện tích.
\end{itemize}
}
\end{tcolorbox}
\end{minipage}

\vspace{-2.5mm}
% --- KHỐI 3: SƠ ĐỒ TƯ DUY TIKZ ---
\noindent
\begin{tcolorbox}[
  enhanced,
  colback=orange!4!white,
  colframe=orange!80!black,
  arc=2mm, boxrule=0.8pt,
  left=2mm, right=2mm, top=1.5mm, bottom=1.5mm,
  title={\footnotesize\bfseries\color{white}\faSitemap\ SƠ ĐỒ TƯ DUY: CÔNG THỨC DIỆN TÍCH CHỮ NHẬT VÀ VUÔNG}
]
\centering
\begin{tikzpicture}[
  box/.style={draw, rounded corners=2pt, align=center, font=\scriptsize\bfseries, inner sep=2.5pt, minimum height=0.75cm},
  arr/.style={-{Stealth[length=2mm]}, thick}
]
  \node[box, fill=purple!15, draw=purple!80!black, text width=2.8cm] (r1) at (0, 0) {HÌNH CHỮ NHẬT\\Dài $a$, Rộng $b$\\$\mathbf{S = a \times b}$};
  \node[box, fill=teal!15, draw=teal!80!black, text width=2.8cm] (r2) at (4.0, 0) {BÀI TOÁN NGƯỢC\\$a = S : b$\\$b = S : a$};
  \node[box, fill=blue!12, draw=blue!80!black, text width=2.8cm] (r3) at (8.0, 0) {HÌNH VUÔNG\\Cạnh $a$\\$\mathbf{S = a \times a}$};
  \node[box, fill=orange!15, draw=orange!80!black, text width=2.4cm] (r4) at (11.2, 0) {LƯU Ý ĐƠN VỊ\\Cùng đơn vị đo\\Kết quả là $\text{cm}^2$};

  \draw[arr, purple!80!black] (r1) -- (r2);
  \draw[arr, blue!80!black] (r2) -- (r3);
  \draw[arr, orange!80!black] (r3) -- (r4);
\end{tikzpicture}
\end{tcolorbox}

\vspace{-2.5mm}
% --- KHỐI 4: VÍ DỤ & MẸO TRÁNH BẪY ---
\noindent
\begin{minipage}[t]{0.49\textwidth}\vspace{0pt}
\begin{tcolorbox}[
  enhanced,
  colback=green!4!white,
  colframe=green!60!black,
  arc=2mm, boxrule=0.8pt,
  left=2mm, right=2mm, top=1.8mm, bottom=1.8mm,
  title={\footnotesize\bfseries\color{white}\faCheckCircle\ Ví dụ thực hành chuẩn mực}
]
{\footnotesize
\textbf{1. Tính diện tích hình chữ nhật:}\\
Hình chữ nhật có chiều dài $12\text{ cm}$, chiều rộng $5\text{ cm}$.\\
Diện tích hình chữ nhật là: $12 \times 5 = \mathbf{60\text{ cm}^2}$.\\
\textbf{2. Tính diện tích hình vuông:}\\
Hình vuông có chu vi là $24\text{ cm}$.\\
Độ dài cạnh hình vuông: $24 : 4 = 6\text{ (cm)}$.\\
Diện tích hình vuông: $6 \times 6 = \mathbf{36\text{ cm}^2}$.
}
\end{tcolorbox}
\end{minipage}
\hfill
\begin{minipage}[t]{0.49\textwidth}\vspace{0pt}
\begin{tcolorbox}[
  enhanced,
  colback=red!4!white,
  colframe=red!70!black,
  arc=2mm, boxrule=0.8pt,
  left=2mm, right=2mm, top=1.8mm, bottom=1.8mm,
  title={\footnotesize\bfseries\color{white}\faExclamationTriangle\ Mẹo tránh bẫy trắc nghiệm}
]
{\footnotesize
\begin{itemize}\setlength{\itemsep}{1.5pt}\setlength{\leftskip}{-2.5mm}
  \item \textbf{Bẫy nhầm $S = a \times a$ với $P = a \times 4$:}\\
  Rất nhiều bạn tính diện tích hình vuông cạnh $5\text{ cm}$ lấy $5 \times 4 = 20$ (đây là chu vi)! Đúng phải là: $5 \times 5 = \mathbf{25\text{ cm}^2}$.
  \item \textbf{Bẫy quên đổi đơn vị đo:}\\
  Dài $1\text{ dm}$, rộng $8\text{ cm}$ $\implies$ phải đổi $1\text{ dm} = 10\text{ cm}$ rồi mới tính $10 \times 8 = \mathbf{80\text{ cm}^2}$.
\end{itemize}
}
\end{tcolorbox}
\end{minipage}


\end{document}
